\scanpage{7}
\mspage{( 2 )}

(ii) En caractéristique $p>0$ (pour les corps résiduels) : si $\mathrm B$ est quasi-séparable, $\mathrm L_b$ est relativement parfait (\textsuperscript{2}) sur $\mathrm L_a$ et
\[
\mathfrak n=(\mathfrak n\cap\mathrm A[\mathrm B^p])\,\mathrm B+\mathfrak n^2,\tag*{(S)}
\]
la réciproque étant vraie si l’on suppose en outre $\unsure{\mathrm D^{\mathrm A}|\mathrm D_{\mathrm A}}(\mathrm B)$ de type fini. Si $\mathrm B$ est nœthérien, (S) se simplifie encore.

\textit{Globalement :}

(i) Soit $\mathrm A$ un anneau nœthérien quasi-séparable sur un corps $\mathrm K$ :

--- en caractéristique $0$, $\mathrm A$ est composé direct d’une famille finie d’extensions algébriques de $\mathrm K$;

--- en caractéristique $p>0$, si pour tout idéal maximal de $\mathrm A$ le corps résiduel est une extension séparable (resp. de type fini) de $\mathrm K$, $\mathrm A$ est composé direct d’une famille finie d’extensions séparables et quasi-séparables (resp. algébriques séparables) de $\mathrm K$.

(ii) Soit $\mathrm A$ une algèbre quasi-séparable sur un corps $\mathrm K$ :

--- en caractéristique $0$, tout idéal premier de $\mathrm A$ est maximal et tous les corps résiduels des localisés de $\mathrm A$ sont des extensions algébriques de $\mathrm K$;

--- en caractéristique $p>0$, si pour tout idéal premier $\mathfrak g$ de $\mathrm A$ le corps $\mathrm A\mathfrak g/\mathfrak g\mathrm A\mathfrak g$ est de type fini sur $\mathrm K$, alors tout idéal premier de $\mathrm A$ est maximal et tous les corps résiduels sont des extensions algébriques séparables de $\mathrm K$;

--- Sous l’une ou l’autre de ces hypothèses, le quotient de $\mathrm A$ par sa nilracine est un anneau absolument plat (\textsuperscript{3}).

2. \textsc{Différentielles et anneaux locaux réguliers.} --- On se place dans le cadre local, en supposant même $\mathrm A$ et $\mathrm B$ intègres, de corps des fractions $\mathrm K_a$ et $\mathrm K_b$, et les cardinaux $\{\mathrm L_b:\mathrm L_a\}$, $\{\mathrm K_b:\mathrm K_a\}$, $\operatorname{rg}\mathfrak n$ finis.

\textsc{Théorème.} --- \textit{Si $\mathrm D_{\mathrm A}(\mathrm B)$ est de type fini (en particulier si $\mathrm B$ est de type fini sur $\mathrm A$) il y a équivalence entre :}

(i) \textit{$\mathrm D_{\mathrm A}(\mathrm B)$ est libre;}

(ii) \textit{a. Dans le cas où $\mathrm L_b$ est séparable sur $\mathrm L_a$ :}
\[
\dim_{\mathrm L_b}\frac{\mathfrak n}{\mathfrak n^2+\mathfrak m\mathrm B}=\{\mathrm K_b:\mathrm K_a\}-\{\mathrm L_b:\mathrm L_a\};
\]

\textit{b. Sinon, $p>0$ désignant la caractéristique de $\mathrm L_b$ et $\mathrm L_a$ :}
\[
\dim_{\mathrm L_b}\frac{\mathfrak n}{\mathfrak n^2+(\mathfrak n\cap\mathrm A[\mathrm B^p])\,\mathrm B}=\{\mathrm K_b:\mathrm K_a\}-\{\mathrm L_b:\mathrm L_a\}.
\]

Notons que l’hypothèse « $\mathrm D_{\mathrm A}(\mathrm B)$ de type fini » entraîne que les cardinaux $\{\mathrm K_b:\mathrm K_a\}$ et $\{\mathrm L_b:\mathrm L_a\}$ sont finis. De plus, l’égalité (ii) \textit{a} s’écrit :

--- Si $\mathrm A$ est un corps : $\operatorname{rang}\mathfrak n=\{\mathrm K_b:\mathrm A\}-\{\mathrm L_b:\mathrm A\}$,

--- Si $\mathrm B$ est non ramifié sur $\mathrm A$ : $\{\mathrm K_b\mid\mathrm K_a\}=\{\mathrm L_b\mid\mathrm L_a\}$.
% archive sheet number: 6
